\scanpage{31}
% archive sheet number: 30
% Typescript is upside down in the supplied scan; transcribed in its reading direction.
\mspage{21}
n\textsuperscript{o} 246

% Printed lines 1--6.
donc, en vertu de la relation (14), $x_{\handwritten{31}}\handwritten{<}0$, donc ($\unsure{\Delta|\handwritten{\Delta}}_{\handwritten{0}},\Delta_{\handwritten{1}},\Delta_{\handwritten{3}}$) est directe. Soit donc $x_{\handwritten{01}}\handwritten{>}0$ ; si $x_{\handwritten{30}}\handwritten{<}0$, on a $x_{\handwritten{02}}\handwritten{>}0$, $x_{\handwritten{23}}\handwritten{>}0$ (parce que la seconde suite est directe), puis $x_{\handwritten{12}}\handwritten{>}0$ (parce que la première est directe), d'où $x_{\handwritten{13}}\handwritten{>}0$, d'après (14), et ($\Delta_{\handwritten{0}},\Delta_{\handwritten{1}},\Delta_{\handwritten{3}},$) est directe. Il en est de même, enfin, si $x_{\handwritten{01}}\handwritten{>}0$, $x_{\handwritten{30}}\handwritten{>}0$. les autres axiomes de la relation d'ordre se vérifient trivialement.

% Printed lines 7--10; underlining does not continue beneath the formula blanks.
Corollaire. - \uline{Soient} $\Delta_i$ ($0\handwritten{\leq}i\handwritten{\leq3}$) \uline{quatre demi-droites telles que la suite} $\Delta_{\handwritten{0}},\Delta_{\handwritten{1}},\Delta_{\handwritten{2}}$ \uline{soit directe}. \uline{Alors}, \uline{pour que la suite} $\Delta_{\handwritten{1}},\Delta_{\handwritten{3}},\Delta_{\handwritten{2}}$ \uline{soit directe}, \uline{il faut et il suffit que les deux suites} $\Delta_{\handwritten{0}},\Delta_{\handwritten{1}},\Delta_{\handwritten{3}}$ \uline{et} $\Delta_{\handwritten{0}},\Delta_{\handwritten{3}},\Delta_{\handwritten{2}}$ \uline{le soient}.

% Printed lines 11--16.
Pour abréger, notons (ijk) l'assertion ``\unsure{lasuite|la suite} $\Delta_i,\Delta_j,\Delta_k$ est directe''. D'après la prop. 6 (132) et (120) entraîne (130) ; et (201) et (213) entraîne (203) ; donc le condition du corollaire est nécessaire. Réciproquement, supposons (012), (013) et (032) vérifiées ; si (123) était vérifiée, alors on aurait (310), car (312) et (320) entraîne (310) d'après la prop. 6. Cette contradiction achève la démonstration.

% Printed lines 17--26; final sentence continues beyond the sheet.
Si $\Delta_{\handwritten{1}},\Delta_{\handwritten{2}}$ sont deux demi-droites distinctes, on voit immédiatement (en prenant une base formée d'un vecteur de $\Delta_{\handwritten{1}}$ et d'un vecteur de $\Delta_{\handwritten{2}}$) qu'il existe une demi-droite $\Delta_{\handwritten{0}}$ telle que la suite $\Delta_{\handwritten{0}},\Delta_{\handwritten{1}},\Delta_{\handwritten{2}}$ soit directe. Il résulte alors de la prop. 6 et de son corollaire que, pour la relation d'ordre déterminée par $\Delta_{\handwritten{0}}$ comme dans la prop. 6, les intervalles ouvert et fermé $\handwritten{]}\Delta_{\handwritten{1}},\Delta_{\handwritten{2}}\handwritten{[}$ et $\handwritten{[}\Delta_{\handwritten{1}},\Delta_{\handwritten{2}}\handwritten{]}$ sont indépendants du choix de $\Delta_{\handwritten{0}}$. La réunion des demi-droites d'un tel intervalle s'appelle un \uline{secteur angulaire}. Autrement dit, on appellera \uline{secteur angulaire ouvert} $\handwritten{]}\Delta_{\handwritten{1}},\Delta_{\handwritten{2}}\handwritten{[}$ la réunion des demi-droites $\Delta$ telles que la suite $\Delta_{\handwritten{1}},\Delta,\Delta_{\handwritten{2}}$ soit directe, et \uline{secteur angulaire fermé} $\handwritten{[}\Delta_{\handwritten{1}},\Delta_{\handwritten{2}}\handwritten{]}$
